\chapter*{Abstract}
\addcontentsline{toc}{chapter}{Abstract}

Positioning systems are extensively used for repetitive tasks and as part of complex robotic systems. In harsh environments like the CERN accelerator complex, integrated commercial electronics frequently fail due to high radiation levels. However, Commercial Off-The-Shelf (COTS) components are widely utilized across CERN's electronics infrastructure, as true radiation-hardened solutions are financially unviable due to astronomical costs. Consequently, these COTS components must undergo a strict Radiation Hardness Assurance (RHA) procedure. This work details the development and evaluation of a prototype radiation-tolerant motor control system built from discrete, radiation-tested COTS components. The results lay the foundation for future design iterations and subsequent integration into robotic and positioning platforms within the CERN accelerator complex.

\vspace{1em}
\textbf{Keywords:} CERN, radiation testing, radiation tolerant, test bench, BLDC, FOC, embedded systems.

\chapter*{Περίληψη}
\addcontentsline{toc}{chapter}{Περίληψη}

\begin{greek}
Συστήματα ηλεκτροκίνησης χρησιμοποιούνται ευρέως για τον αυτοματισμό επαναλαμβανόμενων εργασιών, καθώς και ως τμήματα ρομποτικών συστημάτων. Σε περιβάλλοντα με εξαιρετικά αντίξοες συνθήκες, όπως αυτό του συγκροτήματος επιταχυντών του CERN, τα εμπορικά ηλεκτρονικά συστήματα συχνά αποτυγχάνουν εξαιτίας των πολύ υψηλών επιπέδων ακτινοβολίας. Ωστόσο, μεμονωμένα εξαρτήματα εμπορικού τύπου (COTS) χρησιμοποιούνται στην πλειοψηφία των συστημάτων του CERN σε σημεία όπου η ακτινοβολία είναι χαμηλότερη, καθώς τα radiation-hardened ηλεκτρονικά είναι οικονομικά ασύμφορα. Επομένως, αναπτύσσονται μεγαλύτερα συστήματα που αποτελούνται από COTS ηλεκτρονικά, τα οποία δοκιμάζονται το καθένα ξεχωριστά σε περιβάλλοντα ακτινοβολίας και υποβάλλονται σε αυστηρές διαδικασίες Radiation Hardness Assurance (RHA) προτού ενσωματωθούν. Η παρούσα εργασία παρουσιάζει αναλυτικά τη διαδικασία ανάπτυξης και αξιολόγησης ενός τέτοιου συστήματος ελέγχου κινητήρα, με στόχο τον χαρακτηρισμό του ως επαρκώς ανθεκτικό σε περιβάλλοντα υψηλής ακτινοβολίας. Έπειτα, τα αποτελέσματα θέτουν τις βάσεις για μελλοντικές επεκτάσεις και ενδεχόμενη εγκατάσταση σε ρομποτικά συστήματα εντός των επιταχυντών.

\vspace{1em}
\textbf{Λέξεις-κλειδιά:} CERN, δοκιμές ακτινοβολίας, ανθεκτικότητα σε ακτινοβολία, BLDC, FOC, ενσωματωμένα συστήματα.
\end{greek}

\clearpage